\chapter{Requirements, Specifications, and Evaluation Protocol}
\label{chap:requirements-specifications}

This chapter translates the problem into verifiable expectations. Requirements should be clear enough to guide design and later validation.

\section{Non-Functional Requirements}
Describe quality expectations such as reliability, usability, maintainability, performance, security, scalability, portability, or traceability.

\begin{table}[H]
    \centering
    \caption{Non-functional requirements template}
    \label{tab:non-functional-requirements-template}
    \small
    \begin{tabular}{L{0.14\textwidth}L{0.26\textwidth}L{0.50\textwidth}}
        \toprule
        \textbf{ID} & \textbf{Quality attribute} & \textbf{Expected evidence} \\
        \midrule
        NFR-01 & Replace with a quality attribute. & State how this quality will be checked or observed. \\
        \tabrowsep
        NFR-02 & Replace with a quality attribute. & State how this quality will be checked or observed. \\
        \bottomrule
    \end{tabular}
\end{table}

Avoid vague requirements such as ``the system must be user friendly'' without measurable or observable evidence.

\section{Functional Requirements}
List the services, behaviors, or operations the system must provide. Use active, testable wording.

\begin{table}[H]
    \centering
    \caption{Functional requirements template}
    \label{tab:functional-requirements-template}
    \small
    \begin{tabular}{L{0.14\textwidth}L{0.76\textwidth}}
        \toprule
        \textbf{ID} & \textbf{Requirement} \\
        \midrule
        FR-01 & The system shall replace this sentence with a specific, verifiable behavior. \\
        \tabrowsep
        FR-02 & The system shall replace this sentence with a specific, verifiable behavior. \\
        \bottomrule
    \end{tabular}
\end{table}

Avoid mixing implementation choices into requirements unless they are mandatory constraints imposed by the context.

\section{Technical Constraints}
State constraints related to tools, data, standards, platforms, deadlines, deployment, hardware, software, access rights, or institutional rules.

Avoid presenting personal preferences as constraints. Explain why each constraint exists and how it affects the design.

\section{Success Criteria}
Define what counts as a successful project. Criteria should connect to the objectives and be checked later in the validation chapter.

\section{Evaluation Protocol}
Explain how each requirement will be evaluated. Name the tests, experiments, reviews, demonstrations, datasets, metrics, or acceptance rules that will be used.

\placeholderfigure{evaluation-protocol.pdf}{Evaluation protocol placeholder}{Insert a diagram that links objectives, requirements, tests, and evidence.}

Avoid treating demonstration screenshots as proof unless they are connected to explicit acceptance criteria.

\section{Chapter Summary}
Close by explaining how the requirement baseline guides the methodology and design chapters.